% !TeX root = ../drafts/a-ring-switching.tex
\section{Ring switching}\label{annex:rs}

Ring switching was introduced in \cite{DP24}. We present a variant tailored to our use case: transforming a PCS for $\K$ into a PCS for $\Ftwo$ (transforming a $\K$-valued multilinear PCS, with $\nkeccak$ variables and opening claims in $\E$, into a commitment protocol to 64 $\Ftwo$-valued multilinears, each with $\nkeccak$ variables and sharing opening claims in $\E$).

\subsection{Setting}\label{rs:setting}

The two fields are those of \S\ref{sec:vm}: $\K=\F_2[x]/(x^{64}+x^4+x^3+x+1)$ inside $\E=\K[y]/(y^2+xy+1)$. As $\F_2$-vector spaces, $\K$ has basis $x^0,\dots,x^{63}$, and $\E$ has the $128$ products
\[
e_w=x^iy^j\in\E,\qquad w=64j+i,\quad 0\le i<64,\quad 0\le j<2.
\]

The witness is $64$ bit-valued multilinears $Q_0,\dots,Q_{63}:\cube{\nkeccak}\to\{0,1\}$. We commit to them packed, one $\K$ element per point:
\[
\qkeccak:\cube{\nkeccak}\to\K,\qquad \qkeccak(u)=\sum_{i=0}^{63}Q_i(u)\,x^i.
\]
The commitment scheme proves statements of one shape. Given a target $T\in\E$ and an MLE-friendly weight $W:\cube{\nkeccak}\to\E$ (Definition~\ref{def:ipcs}), an opening proves
\[
\sum_{u\in\cube{\nkeccak}}W(u)\,\qkeccak(u)=T.
\]
The weight $W$ and target $T$ are $\E$-valued, the committed values are $\K$-valued.

What arrives is not of that shape. It is $64$ claimed values $s_i\in\E$, one per bit polynomial, at a point $r\in\E^{\nkeccak}$ they all share.
\[
s_i \qeq \mle{Q_i}(r)=\sum_{u\in\cube{\nkeccak}}\eq(r,u)\,Q_i(u),\qquad 0\le i<64.
\]

Ring switching~\cite{DP24} turns those $64$ claims into a single statement the opening can prove.

\subsection{The reduction}\label{rs:reduction}

Fix an $\F_2$-linear map $\Phi:\E\to\E$:
\begin{equation}
\begin{split}
\sum_{u\in\cube{\nkeccak}}\Phi\bigl(\eq(r,u)\bigr)\,\qkeccak(u)
&=\sum_{i=0}^{63}x^i\sum_{u\in\cube{\nkeccak}}Q_i(u)\,\Phi\bigl(\eq(r,u)\bigr)\\
&=\sum_{i=0}^{63}x^i\,\Phi\Bigl(\sum_{u\in\cube{\nkeccak}}Q_i(u)\,\eq(r,u)\Bigr)
=\sum_{i=0}^{63}x^i\,\Phi(s_i).
\end{split}
\label{eq:complete}
\end{equation}
The reduction is:
\begin{enumerate}
\item once the $s_i$ are fixed, the verifier draws $\Phi$ at random;
\item it computes the target $T=\sum_{i=0}^{63}x^i\Phi(s_i)$;
\item the commitment scheme opens $\qkeccak$ against the weight $W:\cube{\nkeccak}\to\E$, $W(u)=\Phi\bigl(\eq(r,u)\bigr)$, with target $T$; if that opening accepts, the verifier accepts the $64$ claims.
\end{enumerate}
Everything now turns on which $\Phi$ we can afford to draw.

\subsection{Which maps are sound}\label{rs:sound-maps}

\paragraph{The error, transposed.} Suppose the $s_i$ are wrong. The errors $\delta_i=s_i-\mle{Q_i}(r)\in\E$, for $0\le i<64$, are not all zero, and they are fixed before $\Phi$ is drawn. By \eqref{eq:complete} the target the verifier computes differs from the honest one by $\sum_{i=0}^{63}x^i\Phi(\delta_i)$, so the opening succeeds only if
\[
\sum_{i=0}^{63}x^i\,\Phi(\delta_i)=0.
\]
Expand each error in the $\F_2$-basis of $\E$, say $\delta_i=\sum_{w=0}^{127}\delta_{i,w}e_w$ with $\delta_{i,w}\in\F_2$, and exchange the two sums:
\begin{equation}
\sum_{i=0}^{63}x^i\,\Phi(\delta_i)
=\sum_{i=0}^{63}x^i\sum_{w=0}^{127}\delta_{i,w}\,\Phi(e_w)
=\sum_{w=0}^{127}\Phi(e_w)\,\Delta_w,
\qquad
\Delta_w:=\sum_{i=0}^{63}\delta_{i,w}x^i\in\K.
\label{eq:transpose}
\end{equation}
The $x^i$ are an $\F_2$-basis of $\K$, so some $\Delta_w\neq0$. The cheating prover succeeds if $\sum_{w=0}^{127}\Phi(e_w)\,\Delta_w=0$. We must thus draw $\Phi$ to make this event unlikely.

\paragraph{Linear maps are Frobenius sums.} Squaring is additive in characteristic $2$, so
\[
\Phi(a)=\sum_{k=0}^{127}c_k\,a^{2^k},\qquad a\in\E,\quad c_k\in\E,
\]
is $\F_2$-linear for any choice of coefficients. Every $\F_2$-linear map $\E\to\E$ arises this way, and exactly once. Distinct tuples of $(c_k)$ give distinct maps: their difference would be a nonzero polynomial of degree at most $2^{127}$ with $2^{128}$ roots. There are $|\E|^{128}$ tuples and as many maps. Call $\supp=\{k:c_k\neq0\}\subseteq\{0,\dots,127\}$ the \emph{support} of $\Phi$. Substituting into the check above,
\begin{equation}
\sum_{w=0}^{127}\Phi(e_w)\,\Delta_w=\sum_{w=0}^{127}\sum_{k\in\supp}c_k\,e_w^{2^k}\,\Delta_w=\sum_{k\in\supp}c_k\,\pair_k,
\qquad
\pair_k:=\sum_{w=0}^{127}e_w^{2^k}\,\Delta_w\in\E.
\label{eq:pairing}
\end{equation}
The $\pair_k$ depend only on the error, the $c_k$ only on the randomness. So we need the map $\K^{128}\to\E^{|\supp|}$, $\Delta\mapsto(\pair_k)_{k\in\supp}$, to be injective: a nonzero error has to leave some $\pair_k\neq0$ for the coefficients to act on.

\paragraph{Sixty-four terms.} That map is $\K$-linear, since the $\Delta_w$ enter linearly and the Frobenius touches only the constants $e_w$. It sends $\K^{128}$, of $\K$-dimension $128$, into $\E^{|\supp|}$, of $\K$-dimension $2|\supp|$, so injectivity needs $|\supp|\ge\lceil128/2\rceil=64$. Take $\supp=\{0,\dots,63\}$ and suppose $\pair_k=0$ for all $k<64$. Two steps force $\Delta=0$.

First, split $\pair_k$ over $\K$. Write $w=64j+i$, so that $e_w^{2^k}=x^{i2^k}\,\bigl(y^{2^k}\bigr)^{j}$ and
\[
\pair_k=\pair_{0,k}+y^{2^k}\,\pair_{1,k},
\qquad
\pair_{j,k}:=\sum_{i=0}^{63}\Delta_{64j+i}\,x^{i2^k}\in\K,\qquad 0\le j<2.
\]
Squaring is a bijection of $\K$ onto itself, so $y^{2^k}\in\K$ would force $y\in\K$, hence $\E=\K[y]=\K$. By Lemma~\ref{lem:kbasis} applied to $y^{2^k}$, the set $\bigl\{1,y^{2^k}\bigr\}$ is a $\K$-basis of $\E$, and $\pair_k=0$ forces $\pair_{0,k}=\pair_{1,k}=0$.

Second, read the $64$ equations $\pair_{j,k}=0$, at fixed $j$, as a single polynomial. Put
\[
D_j(v)=\sum_{i=0}^{63}\Delta_{64j+i}\,v^i\in\K[v],
\qquad\text{so that}\qquad
\pair_{j,k}=D_j\bigl(x^{2^k}\bigr).
\]
Then $D_j$ has degree at most $63$ and vanishes at the $64$ points $x,x^2,x^4,\dots,x^{2^{63}}$. Those are distinct, because $x$ has degree $64$ over $\F_2$ and so has $64$ distinct conjugates. A nonzero polynomial of degree at most $63$ cannot have $64$ roots, so $D_j=0$ and every $\Delta_{64j+i}$ vanishes. A nonzero error therefore leaves some $\pair_k\neq0$.

The count is tight on both sides: $64$ terms give $64$ equations $\pair_k=0$ over $\E$, which is $128$ equations over $\K$, against exactly $128$ unknowns $\Delta_w\in\K$.

\subsection{The map we draw}\label{rs:map}

We need a random $\Phi$ of support $\{0,\dots,63\}$, cheap to apply. Instead of drawing its $64$ coefficients independently, we derive them from sixteen challenges, in two stages of eight.

Draw $f_0,\dots,f_7,g_0,\dots,g_7\leftarrow\E$, in this order, once the $s_i$ are fixed, and put, for $a\in\E$,
\begin{equation}
\Psi(a)=\sum_{b<8}g_b\,a^{2^b},
\qquad
\Phi(a)=\sum_{p<8}f_p\,\Psi(a)^{2^{8p}}.
\label{eq:phi}
\end{equation}
Both stages are $\F_2$-linear, so $\Phi$ is. Squaring is additive, so $\Psi(a)^{2^{8p}}=\sum_bg_b^{2^{8p}}a^{2^{8p+b}}$, and with $k=8p+b$ the composition expands to
\begin{equation}
\Phi(a)=\sum_{k=0}^{63}c_k\,a^{2^k},
\qquad
c_{8p+b}=f_p\,g_b^{\,2^{8p}}\qquad(p,b<8).
\label{eq:ck}
\end{equation}
Each $k$ from $0$ to $63$ occurs exactly once, one term for every Frobenius power the previous section requires.

\paragraph{Soundness.} Fix a nonzero error before the challenges are drawn. The previous section gives some $\pair_{8p+b}\neq0$. Group the check by $p$, writing $a^{2^{-8p}}$ for the inverse of the Frobenius $a\mapsto a^{2^{8p}}$, a bijection of $\E$:
\[
\sum_{k=0}^{63}c_k\,\pair_k=\sum_{p<8}f_p\,M_p,
\qquad
M_p=\sum_{b<8}g_b^{\,2^{8p}}\,\pair_{8p+b}=\Bigl(\sum_{b<8}g_b\,\pair_{8p+b}^{\,2^{-8p}}\Bigr)^{2^{8p}}.
\]
For that $p$ the inner sum is a nonzero linear form in $g_0,\dots,g_7$, so $M_p=0$ with probability $1/|\E|$; when $M_p\neq0$, the outer sum is a nonzero linear form in $f_0,\dots,f_7$, drawn independently, and vanishes with probability $1/|\E|$. The check passes with probability at most $2/|\E|$ for one committed polynomial. That the $c_k$ reach degree $2^{56}$ in the $g_b$ does not matter: each stage is linear in its own challenges, and the Frobenius only permutes $\E$. Against the list of the committed oracle (Annex~\ref{annex:pcs}), a union bound pays this once per member; the protocol draws the sixteen challenges with the opening's batching challenge, after one grind of $\grind{batch}$ bits that brings the whole row to $2^{-123}$ (\S\ref{sec:soundness-outer}).

One stage of sixty-four terms would take sixty-four challenges; two stages of eight take sixteen, and putting the large Frobenius powers outside lets the target share them across the $64$ slices (\S\ref{rs:cost}).

\subsection{Evaluating the weight}\label{rs:weight}

An opening reduces $\sum_{u\in\cube{\nkeccak}}W(u)\,\qkeccak(u)=T$ to a claim about $\qkeccak$ at a random point $r'\in\E^{\nkeccak}$, and the verifier must evaluate the multilinear extension of $W$ at $r'$, written $\mle{W}(r')$, by itself. One fact turns that into a short formula.

\emph{Fact (Frobenius passes through $\eq$ on the boolean hypercube).} Squaring respects sums and products and fixes bits, so for $u\in\cube{\nkeccak}$, by Definition~\ref{def:eq},
\[
\eq(r,u)^{2^k}=\prod_{n=1}^{\nkeccak}(1+r_n+u_n)^{2^k}=\prod_{n=1}^{\nkeccak}\bigl(1+r_n^{2^k}+u_n\bigr)=\eq\bigl(r^{2^k},u\bigr),
\qquad r^{2^k}:=\bigl(r_1^{2^k},\dots,r_{\nkeccak}^{2^k}\bigr).
\]

Expand $W$ at $r'$ with \eqref{eq:ck} and this fact:
\[
\begin{aligned}
\mle{W}(r')&=\sum_{u\in\cube{\nkeccak}}W(u)\,\eq(r',u)
=\sum_{u\in\cube{\nkeccak}}\Phi\bigl(\eq(r,u)\bigr)\,\eq(r',u)\\
&=\sum_{k=0}^{63}c_k\sum_{u\in\cube{\nkeccak}}\eq\bigl(r^{2^k},u\bigr)\,\eq(r',u).
\end{aligned}
\]
The inner sum is $\eq(r^{2^k},r')$, since $\eq(u,\cdot\,)$ is its own multilinear extension for every $u$. By Definition~\ref{def:eq} again,
\[
\boxed{\ \
\mle{W}(r') \;=\; \sum_{k=0}^{63}\, c_k \prod_{n=1}^{\nkeccak}
\bigl( 1 + r_n^{2^k} + r'_n \bigr).
\ \ }
\]

\subsection{Cost}\label{rs:cost}

The target need not apply $\Phi$ to each slice. Put $z_p=x^{2^{-8p}}$, a constant of $\K$, equal to $x^{2^{64-8p}}$ since the Frobenius of $\K$ has order $64$. Then $x^i\,\Psi(a)^{2^{8p}}=\bigl(z_p^i\,\Psi(a)\bigr)^{2^{8p}}$, and
\[
T=\sum_{i=0}^{63}x^i\,\Phi(s_i)=\sum_{p<8}f_p\Bigl(\sum_{i=0}^{63}z_p^i\,\Psi(s_i)\Bigr)^{2^{8p}}.
\]
One application of $\Psi$ takes $7$ squarings and $8$ multiplications. The verifier applies it to each slice, takes the eight inner sums by Horner in $z_p$, raises the $p$-th to $2^{8p}$, $224$ squarings in all, and multiplies by $f_p$. The coefficients of \eqref{eq:ck} take $56$ squarings per $g_b$ and one multiplication per $c_k$.

\begin{center}
\begin{tabular}{lrr}
& squarings & multiplications\\[2pt]
\hline
one application of $\Psi$ & $7$ & $8$\\
the target $T$, $64$ applications and $8$ sums & $672$ & $520$\\
the coefficients $c_k$, once & $448$ & $64$\\
the weight $\mle{W}(r')$ & $63\nkeccak$ & $64\nkeccak$\\
\end{tabular}
\end{center}

Everything above is in $\E$. The $504$ Horner steps in the target are left out: each multiplies by a constant $z_p$ of $\K$, which is cheaper.

\subsection{Credits}\label{rs:credits}

The construction above rests on generalizing ring switching~\cite{DP24} beyond the binary-tower setting, an idea of Lev Soukhanov, \textit{[[alloc]init]}. The idea of drawing $\Phi$ as a composition of cheap stages is from~\cite{Flock26}.
